\section{What \adapt{} Offers Relative to Prior Work}
\label{sec:compare}

Table~\ref{tab:gap} shows that \adapt{} covers more stages of the image-to-mission chain than any single group of reviewed works. This breadth is its main contribution, and it comes with explicit trade-offs.

\emph{End-to-end executable output.} Flood-mapping methods \cite{twele2016sentinel,gebrehiwot2019deep,simantiris2024unsupervised} end at a mask, and relief-routing models \cite{rabta2018drone,dong2026path} end at a schedule over given nodes. \adapt{} produces a file that MAVLink ground stations are designed to load, containing explicit delivery actions. The trade-off is that each stage is simpler than its specialised counterpart: colour thresholding instead of learned or spectral segmentation, a heuristic spread rule instead of a hydrodynamic or CA inundation model \cite{jamali2019cellular}, a nearest-neighbour heuristic instead of exact or metaheuristic optimisation, and boundary geometry instead of site-safety scoring \cite{patterson2014timely}.

\emph{Low data and compute requirements.} The pipeline needs no training data, no digital elevation model and no GPU. It runs on a laptop in under 40~s for the evaluated maps, with planning taking almost all of that time; on these inputs a static A* planner would reduce planning to under 0.3~s (Table~\ref{tab:baselines}). The price is that nothing in the pipeline is calibrated against physical flood behaviour or delivery safety.

\emph{Weather-informed obstacles.} Unlike static-map planners \cite{aggarwal2020path}, \adapt{} inflates obstacles according to current precipitation and wind. Our ablation shows that this reduces the route length over predicted flooding by 47\% and 76\% on the two dot maps, but disconnects drop points that the planner then reaches by straight, unchecked legs. Prior CA flood models are physically richer but are not coupled to planning \cite{ghimire2013formulation,guidolin2016weighted}.

\emph{Verified geospatial and mission layers.} We found no reviewed image-to-UAV-mission work that reports independent verification of its coordinate conversion and mission encoding. \adapt{} provides an independent geodetic oracle, a tested validator and mutation evidence, and states the conversion's validity domain. This strengthens confidence in the software, but it does not substitute for geographic validation of the inputs or for flight testing. Because code, evaluation scripts and results are released together, any component (for example a learned segmenter or a CA flood model) can be replaced and re-evaluated in place.
